%% RQ4 results (Phase 7). Every number is a macro from tables/numbers.tex.
Every result in this subsection comes from the setting in which the arms differ only in their exporter: the common population of conversations present in every arm, the schedule ground truth and the conversation-grouped split. The analysis plan for RQ4, with the rule for a large residual and every sensitivity variant, was committed to the repository before any RQ4 number was computed.

\paragraph{H4 is not supported, and the test cannot decide it.} Table~\ref{tab:rq4-h4} gives Kendall's $W$ of the model ranking across arms. On \cicids{} it is \RQfourWFlow{} at the flow unit and \RQfourWConv{} at the conversation unit. The difference, \RQfourDiff{} (interval \RQfourDiffLo{} to \RQfourDiffHi{}), points away from the hypothesis, but its interval includes zero, so by the rule fixed in advance H4 is not supported and is not reversed either. On UNSW-NB15, with three arms, $W$ is \RQfourWFlowUnsw{} at the flow unit and \RQfourWConvUnsw{} at the conversation unit; the difference, \RQfourDiffUnsw{} (\RQfourDiffLoUnsw{} to \RQfourDiffHiUnsw{}), also includes zero, so H4 is not supported there either. The test has little power on \cicids{} as well. With three models $W$ can take only a few values, and the multilayer perceptron ranks last in \RQtwoMlpLastArms{} arm (Section~\ref{sec:rq2}), so the rankings differ only in whether the random forest or LightGBM comes first. A high $W$ also says that the arms order the models alike, not that their scores agree. A result of ``not supported'' here is therefore not evidence that rankings fail to converge. To give the test more room we added logistic regression and a $k$-nearest-neighbour classifier, tuned and run like the other three (Section~\ref{sec:protocol}). With five models, $W$ on \cicids{} is \RevFiveWFlowCic{} at the flow unit and \RevFiveWConvCic{} at the conversation unit, and the interval of the difference, \RevFiveWDiffLoCic{} to \RevFiveWDiffHiCic{}, lies below zero: by the same rule H4 would be reversed, and the arms rank the models less alike at the conversation unit. On UNSW-NB15 $W$ is \RevFiveWFlowUnsw{} and \RevFiveWConvUnsw{} (difference \RevFiveWDiffLoUnsw{} to \RevFiveWDiffHiUnsw{}), and H4 is again not supported. This grid was added after the primary results, so it does not replace the verdict fixed in advance, but it rules out the reading that a test with more models would have shown rankings converging. The two added models also change H2. Logistic regression has the lowest macro-F1 in all \RevFiveArms{} arms of the two benchmarks, so the models now differ more, and the exporter leads for fewer classes: for \RevFiveAboveCic{} of \RevEligCic{} eligible \cicids{} classes, so H2 still holds there, but for \RevFiveAboveUnsw{} of \RevEligUnsw{} UNSW-NB15 classes, short of a strict majority (\RevFivePerdfAboveUnsw{} per degree of freedom). The exporter share is relative to the models compared, and on UNSW-NB15 the H2 verdict depends on the set of models as well as on CIC-UNSW-NB15.

% Generated by 03_Code/make_numbers.py from 04_Results/phase7/h4.csv. Do not edit by hand.
\begin{sidewaystable}
\caption{H4: Kendall's $W$ of the model ranking across arms at the flow and the conversation unit,
grouped split, common population, mean over seeds. Brackets: 95 percent cluster bootstrap intervals
(1,000 resamples of conversations). H4 predicts a positive difference (conversation minus
flow). Undecided: the interval includes zero, so H4 is not supported.}\label{tab:rq4-h4}
\small
\begin{tabularx}{\linewidth}{@{}lLllll@{}}
\toprule
Benchmark & Setting & $W$, flow & $W$, conversation & Difference & Verdict \\
\midrule
CICIDS2017 & Primary: $T = 15$ s, any rule, schedule labels & 0.886 [0.780, 0.927] & 0.820 [0.763, 0.869] & $-0.065$ [$-0.090$, 0.008] & undecided \\
 & $T = 5$ s & 0.886 [0.788, 0.910] & 0.820 [0.763, 0.861] & $-0.065$ [$-0.090$, 0.016] & undecided \\
 & $T = 30$ s & 0.886 [0.780, 0.910] & 0.820 [0.763, 0.869] & $-0.065$ [$-0.090$, 0.016] & undecided \\
 & $T = 60$ s & 0.886 [0.780, 0.918] & 0.820 [0.763, 0.870] & $-0.065$ [$-0.090$, 0.016] & undecided \\
 & $T = 120$ s & 0.886 [0.788, 0.910] & 0.820 [0.763, 0.869] & $-0.065$ [$-0.090$, 0.016] & undecided \\
 & Majority rule & 0.886 [0.780, 0.927] & 0.820 [0.771, 0.869] & $-0.065$ [$-0.090$, 0.016] & undecided \\
 & Byte-weighted rule & 0.886 [0.780, 0.927] & 0.837 [0.771, 0.878] & $-0.049$ [$-0.090$, 0.016] & undecided \\
 & Episode unit & 0.812 [0.747, 0.837] & 0.751 [0.706, 0.799] & $-0.062$ [$-0.100$, 0.005] & undecided \\
 & Attempted folded into benign (scoring only) & 0.837 [0.771, 0.861] & 0.502 [0.461, 0.551] & $-0.335$ [$-0.355$, $-0.269$] & reversed \\
 & Attempted folded into benign (retrained) & 0.616 [0.347, 0.804] & 0.535 [0.298, 0.739] & $-0.082$ [$-0.159$, 0.041] & undecided \\
 & Each release's own labels (scoring only) & 0.167 [0.151, 0.257] & 0.176 [0.143, 0.249] & 0.008 [$-0.049$, 0.024] & undecided \\
 & Each release's own labels (retrained) & 0.237 [0.131, 0.306] & 0.216 [0.110, 0.290] & $-0.020$ [$-0.078$, 0.041] & undecided \\
\midrule
UNSW-NB15 & Primary: $T = 15$ s, any rule, schedule labels & 0.867 [0.778, 0.956] & 0.822 [0.778, 0.911] & $-0.044$ [$-0.089$, 0.044] & undecided \\
 & $T = 5$ s & 0.867 [0.778, 0.911] & 0.822 [0.778, 0.911] & $-0.044$ [$-0.089$, 0.044] & undecided \\
 & $T = 30$ s & 0.867 [0.778, 0.911] & 0.822 [0.778, 0.911] & $-0.044$ [$-0.089$, 0.044] & undecided \\
 & $T = 60$ s & 0.867 [0.778, 0.911] & 0.822 [0.778, 0.911] & $-0.044$ [$-0.089$, 0.044] & undecided \\
 & $T = 120$ s & 0.867 [0.778, 0.956] & 0.822 [0.778, 0.911] & $-0.044$ [$-0.089$, 0.044] & undecided \\
 & Majority rule & 0.867 [0.778, 0.956] & 0.822 [0.778, 0.911] & $-0.044$ [$-0.133$, 0.044] & undecided \\
 & Byte-weighted rule & 0.867 [0.778, 0.956] & 0.822 [0.778, 0.911] & $-0.044$ [$-0.089$, 0.044] & undecided \\
 & Episode unit & 0.822 [0.778, 0.868] & 0.822 [0.778, 0.867] & 0.000 [$-0.044$, 0.044] & undecided \\
 & Each release's own labels (scoring only) & 0.867 [0.778, 0.956] & 0.867 [0.778, 0.956] & 0.000 [$-0.044$, 0.089] & undecided \\
\bottomrule
\end{tabularx}
\end{sidewaystable}


\paragraph{Definition~\ref{def:invariance} tested directly.} H4 asks whether the arms order the models alike; Definition~\ref{def:invariance} asks whether a metric has the same value whatever the exporter. We test that directly on the common population with two one-sided tests: for each model, class with enough test conversations in every arm, and pair of arms, the per-class F1 of the two arms, averaged over seeds, is called equivalent at margin $\delta$ when the \setnum{90} percent cluster bootstrap interval of its difference lies inside $(-\delta, \delta)$. We report $\delta = \setnum{0.01}$, the margin of the RQ5 claims of no difference, \setnum{0.02} and \setnum{0.05}, a difference most papers would report as an improvement. Over the \RevEqTestsCic{} tests on \cicids{}, \RevEqCicFlowTwo{} percent are equivalent at \setnum{0.02} at the flow unit and \RevEqCicConvTwo{} percent at the conversation unit, and \RevEqCicFlowFive{} and \RevEqCicConvFive{} percent at \setnum{0.05}. For macro-F1 over the fixed classes the share at \setnum{0.05} rises from \RevEqMacCicFlowFive{} to \RevEqMacCicConvFive{} percent. On UNSW-NB15, over \RevEqTestsUnsw{} tests, the shares are much lower: \RevEqUnswFlowFive{} percent at the flow unit and \RevEqUnswConvFive{} percent at the conversation unit at \setnum{0.05}, and no macro-F1 comparison is equivalent at \setnum{0.02} at either unit. So the conversation unit moves both benchmarks towards invariance, but by Definition~\ref{def:invariance} neither metric is exporter invariant on either benchmark: even at the conversation unit, the comparisons not counted above, a minority on \cicids{} and most on UNSW-NB15, cannot be shown to differ by less than \setnum{0.05}.

\paragraph{The conversation unit removes most of the exporter share, but not all of it.} The exporter share of the variance in per-class F1 tells more than the ranking does. From the flow unit to the conversation unit it falls, with an interval that excludes zero, for \RQfourShareFallsCic{} of the \RQfourShareOfCic{} \cicids{} classes with enough test conversations, benign included: SSH-Patator goes from \RQfourSshShareFlow{} to \RQfourSshShareConv{}, DoS GoldenEye from \RQfourGoldenShareFlow{} to \RQfourGoldenShareConv{} and Bot from \RQfourBotShareFlow{} to \RQfourBotShareConv{}. This is what the RQ3 mechanism predicts, since scoring each conversation once removes the weight that tail records add. For \RQfourShareRisesCic{} classes the share rises: DoS Hulk from \RQfourHulkShareFlow{} to \RQfourHulkShareConv{} and FTP-Patator from \RQfourFtpShareFlow{} to \RQfourFtpShareConv{}. On UNSW-NB15 the share falls for \RQfourShareFallsUnsw{} of \RQfourShareOfUnsw{} classes, most for Generic, from \RQfourGenericShareFlow{} to \RQfourGenericShareConv{}. At the conversation unit the exporter leads for only \RQfourUnswConvAbove{} of the \RQfourUnswConvElig{} eligible UNSW-NB15 classes, so on that benchmark H2 holds at the flow unit and not at the conversation unit.

\paragraph{What remains.} We call a residual large when the exporter share at the conversation unit is at least \setnum{0.25} and the arm means of F1 spread by at least \setnum{0.05}. Table~\ref{tab:rq4-residual} lists the \RQfourLargeCic{} \cicids{} classes and \RQfourLargeUnsw{} UNSW-NB15 classes that meet both conditions.

% Generated by 03_Code/make_numbers.py from 04_Results/phase7/residual.csv. Do not edit by hand.
\begin{table}[t]
\caption{Classes with a large residual exporter effect at the conversation unit (grouped split,
common population): exporter share of the variance in per-class F1 at least \setnum{0.25} and a
spread of arm means of F1 of at least \setnum{0.05}. Brackets: 95 percent cluster bootstrap
intervals. Last column: the two arms with the lowest mean F1 over three models and five
seeds.}\label{tab:rq4-residual}
\small
\begin{tabularx}{\linewidth}{@{}lllL@{}}
\toprule
Class & Exporter share & Spread & Lowest arms (mean F1) \\
\midrule
\multicolumn{4}{@{}l}{\emph{\cicids}} \\
FTP-Patator & 0.966 [0.949, 0.973] & 0.085 & CIC original 0.915; BCCC 0.972 \\
DoS Hulk & 0.722 [0.708, 0.734] & 0.128 & NFS-2023-TE 0.872; BCCC 0.945 \\
Bot & 0.559 [0.508, 0.587] & 0.070 & BCCC 0.929; CIC original 0.944 \\
SSH-Patator & 0.301 [0.292, 0.309] & 0.065 & BCCC 0.933; CIC original 0.994 \\
\midrule
\multicolumn{4}{@{}l}{\emph{UNSW-NB15}} \\
Fuzzers & 0.841 [0.810, 0.855] & 0.067 & CIC-UNSW-NB15 0.797; UNSW-NB15 original, duplicates dropped 0.851 \\
Shellcode & 0.788 [0.752, 0.818] & 0.467 & CIC-UNSW-NB15 0.366; UNSW-NB15 original, duplicates dropped 0.658 \\
Exploits & 0.571 [0.525, 0.601] & 0.050 & CIC-UNSW-NB15 0.820; UNSW-NB15 original, duplicates dropped 0.848 \\
Analysis & 0.479 [0.274, 0.587] & 0.162 & NF-UNSW-NB15-v3 0.202; CIC-UNSW-NB15 0.243 \\
Worms & 0.329 [0.234, 0.393] & 0.267 & UNSW-NB15 original, duplicates dropped 0.426; CIC-UNSW-NB15 0.484 \\
DoS & 0.280 [0.241, 0.313] & 0.136 & CIC-UNSW-NB15 0.399; UNSW-NB15 original, duplicates dropped 0.420 \\
Generic & 0.272 [0.244, 0.299] & 0.118 & CIC-UNSW-NB15 0.751; UNSW-NB15 original, duplicates dropped 0.769 \\
\bottomrule
\end{tabularx}
\end{table}


On \cicids{} the residual has one cause, and it is the mirror image of RQ3. For the four classes, recall is at least \RQfourResRecallMin{} in every arm. What falls is precision, and only in the three arms that write tail records: CIC original, BCCC and NFS-2023-TE. In the other four arms precision stays at \RQfourPrecOtherMin{} or above. Almost every false positive in the three tail arms is a benign conversation with two or more records, exactly one of which the model assigned to the attack class: \RQfourFpOneFtp{} of \RQfourFpFtp{} FTP-Patator false positives in CIC original, \RQfourFpOneSsh{} of \RQfourFpSsh{} SSH-Patator false positives in BCCC, \RQfourFpOneBot{} of \RQfourFpBot{} Bot false positives in BCCC, and \RQfourFpOneHulk{} of \RQfourFpHulk{} DoS Hulk false positives in NFS-2023-TE (counts summed over three models and five seeds). BCCC DoS Hulk differs in degree: \RQfourFpMultiHulkBccc{} of its \RQfourFpHulkBccc{} false positives are benign conversations with several records, but only \RQfourFpOneHulkBccc{} have exactly one record flagged. The schedule labels the tail record of an attack conversation as the attack, because a record inherits the label of its conversation. A model trained on these labels learns that a short end-of-connection record belongs to the attack class, and the same kind of record at the end of a benign connection then raises an alert, which the ``any'' rule turns into a false alarm for the whole conversation. In a spot check made by hand from the stored predictions and not part of the scripted pipeline (one model and seed per class), the flagged record in CIC original and BCCC was in most cases the smallest record of its conversation, with zero bytes in the release; we report it as an unscripted observation.

NFS-2023-TE DoS Hulk is only partly explained. Its flagged records are not the smallest of their conversations, and the conversations are ordinary web connections to many external hosts on every capture day, Monday included, which has no attack, so they are not mislabelled attack traffic. A record that TE writes in benign web connections resembles the records TE writes for DoS Hulk; which expiry rule produces it we have not confirmed. This is the same set of benign records as in RQ3 (Section~\ref{sec:rq3}).

On UNSW-NB15 the residual comes mostly from CIC-UNSW-NB15, and it is not the \cicids{} pattern. CIC-UNSW-NB15 has the lowest mean F1 of the three arms in \RQfourUnswCicLowestFone{} of the \RQfourLargeUnsw{} classes and the lowest recall in \RQfourUnswCicLowestRecall{}. In that release \RQfourUnswCicMultiMin{} to \RQfourUnswCicMultiMax{} percent of the conversations of these classes have several records, against at most \RQfourUnswOtherMultiMax{} percent in the other two, and only \RQfourUnswCicFnBenign{} percent of the attack conversations it misses are called benign: the models confuse attack classes with each other rather than with benign traffic. Whether the split records cause these confusions, or the CICFlowMeter features of a split record, we did not test. The original and NF-UNSW-NB15-v3 also differ in their feature sets, the Argus and Bro features against the NetFlow fields, and which features matter we did not test either.

Three checks support the reading for \cicids{}. Under the payload-aware labels of Section~\ref{sec:rq3}, which call records without payload benign in training and test, \RevPayLarge{} \cicids{} classes have a large residual, against \RevSchLarge{} under the schedule, and no attack class spreads by more than \RevPaySpreadMax{} F1 across arms at the conversation unit. $W$ under these labels is \RevPayWFlow{} at the flow unit and \RevPayWConv{} at the conversation unit (difference interval \RevPayWDiffLo{} to \RevPayWDiffHi{}), so H4 is not supported there either. Under the byte-weighted rule, which assigns a conversation the class of its records with the most bytes and so ignores an empty tail, \RQfourBytesLarge{} \cicids{} classes have a large residual and the arm means of the four classes spread by at most \RQfourBytesSpreadMax{}. Leaving out the arm that writes the tails lowers the residual share (point estimates, no intervals): for SSH-Patator from \RQfourSshShareConv{} to \RQfourLoaoSshBccc{} without BCCC, and for DoS Hulk from \RQfourHulkShareConv{} to \RQfourLoaoHulkTe{} without NFS-2023-TE. For FTP-Patator, leaving out CIC original lowers it only to \RQfourLoaoFtpOrig{}, because BCCC writes the same tails.

\paragraph{Sensitivity.} Table~\ref{tab:rq4-sens} states, for each of the \RQfourSensVariants{} \cicids{} variants, which conclusions change. Five conclusions are tracked: H4, H2 at the flow unit, H2 at the conversation unit, the per-class verdict of exporter against model at the conversation unit, and the direction in which the exporter share moves from flow to conversation. H4 changes under \RQfourSensChangedHfour{} variant, H2 at the flow unit under \RQfourSensChangedHtwoFlow{} and H2 at the conversation unit under \RQfourSensChangedHtwoConv{}. Class verdicts change for at most \RQfourSensKfourMax{} of \RQfourClassesOf{} classes, and share directions for at most \RQfourSensKfiveMax{}. On UNSW-NB15, H2 holds at the flow unit and not at the conversation unit in every variant except scoring against each release's own labels, under which the exporter leads at the flow unit for only \RQfourUnswOwnAbove{} of \RQfourUnswOwnElig{} classes and H2 does not hold there either.

% Generated by 03_Code/make_numbers.py from 04_Results/phase7/sensitivity_table.csv. Do not edit by hand.
\begin{sidewaystable}
\caption{What each sensitivity setting changes. H4: the difference in Kendall's $W$, conversation
minus flow, with its interval. H2 at the flow and the conversation unit. Class verdicts: how many
classes change their exporter-against-model verdict at the conversation unit (grouped split, common
population). Share direction: how many classes change the direction in which the exporter share
moves from the flow to the conversation unit. A class whose verdict becomes not estimable is not
counted as a change. Changed: the conclusions that differ from the primary setting.}\label{tab:rq4-sens}
\small
\setlength{\tabcolsep}{4pt}
\begin{tabularx}{\linewidth}{@{}lLlllrrl@{}}
\toprule
 & & & \multicolumn{2}{c}{H2} & Class & Share & \\
\cmidrule(lr){4-5}
Benchmark & Setting & H4 difference & Flow & Conversation & verdicts & direction & Changed \\
\midrule
CICIDS2017 & Primary: $T = 15$ s, any rule, schedule labels & $-0.065$ [$-0.090$, 0.008] & holds & holds & 0 of 12 & 0 of 12 & none \\
 & $T = 5$ s & $-0.065$ [$-0.090$, 0.016] & holds & holds & 1 of 12 & 0 of 12 & none \\
 & $T = 30$ s & $-0.065$ [$-0.090$, 0.016] & holds & holds & 0 of 12 & 0 of 12 & none \\
 & $T = 60$ s & $-0.065$ [$-0.090$, 0.016] & holds & holds & 0 of 12 & 1 of 12 & none \\
 & $T = 120$ s & $-0.065$ [$-0.090$, 0.016] & holds & holds & 0 of 12 & 0 of 12 & none \\
 & Majority rule & $-0.065$ [$-0.090$, 0.016] & holds & holds & 3 of 12 & 2 of 12 & none \\
 & Byte-weighted rule & $-0.049$ [$-0.090$, 0.016] & holds & does not hold & 4 of 12 & 2 of 12 & H2 conv. \\
 & Episode unit & $-0.062$ [$-0.100$, 0.005] & not estimable & not estimable & 0 of 12 & 1 of 12 & none \\
 & Attempted folded into benign (scoring only) & $-0.335$ [$-0.355$, $-0.269$] & holds & holds & 0 of 12 & 1 of 12 & H4 \\
 & Attempted folded into benign (retrained) & $-0.082$ [$-0.159$, 0.041] & holds & holds & 2 of 12 & 2 of 12 & none \\
 & Each release's own labels (scoring only) & 0.008 [$-0.049$, 0.024] & holds & holds & 2 of 12 & 7 of 12 & none \\
 & Each release's own labels (retrained) & $-0.020$ [$-0.078$, 0.041] & holds & holds & 2 of 12 & 9 of 12 & none \\
\midrule
UNSW-NB15 & Primary: $T = 15$ s, any rule, schedule labels & $-0.044$ [$-0.089$, 0.044] & holds & does not hold & 0 of 10 & 0 of 10 & none \\
 & $T = 5$ s & $-0.044$ [$-0.089$, 0.044] & holds & does not hold & 0 of 10 & 0 of 10 & none \\
 & $T = 30$ s & $-0.044$ [$-0.089$, 0.044] & holds & does not hold & 0 of 10 & 0 of 10 & none \\
 & $T = 60$ s & $-0.044$ [$-0.089$, 0.044] & holds & does not hold & 0 of 10 & 0 of 10 & none \\
 & $T = 120$ s & $-0.044$ [$-0.089$, 0.044] & holds & does not hold & 0 of 10 & 0 of 10 & none \\
 & Majority rule & $-0.044$ [$-0.133$, 0.044] & holds & does not hold & 2 of 10 & 5 of 10 & none \\
 & Byte-weighted rule & $-0.044$ [$-0.089$, 0.044] & holds & does not hold & 1 of 10 & 0 of 10 & none \\
 & Episode unit & 0.000 [$-0.044$, 0.044] & holds & does not hold & 1 of 10 & 1 of 10 & none \\
 & Each release's own labels (scoring only) & 0.000 [$-0.044$, 0.089] & does not hold & does not hold & 2 of 10 & 1 of 10 & H2 flow \\
\bottomrule
\end{tabularx}
\end{sidewaystable}


The gap $T$ does not matter between \NsweepCicStableMin{} and \setnum{120} seconds: each of the four alternative values changes at most \RQfourSensTClassMax{} class verdict or share direction, and no hypothesis verdict. The aggregation rule matters for the residual and not for H4. The byte-weighted rule removes the \cicids{} residual, as above, and with it the exporter's lead at the conversation unit, so H2 at that unit no longer holds; majority voting keeps the residual, because most of these conversations have two records and a tie goes to the attack class. An evaluator who scores per conversation therefore has to state the rule.

Folding attempted attacks into benign at test time, with models trained on the primary labels, is the one variant that reverses H4 (difference \RQfourFoldScoreDiff{}, interval \RQfourFoldScoreDiffLo{} to \RQfourFoldScoreDiffHi{}). This variant scores the models as wrong for calling attempts attacks, which they were trained to do. Retrained on the folded labels, the reversal is gone: the difference is \RQfourFoldRetrDiff{} (\RQfourFoldRetrDiffLo{} to \RQfourFoldRetrDiffHi{}), undecided as in the primary setting, and \RQfourFoldRetrainedClasses{} class verdicts change. We read the reversal as a product of scoring against labels the models did not learn.

Scored against each release's own labels instead of the schedule, the arms nearly stop agreeing on the model ranking: $W$ at the flow unit is \RQfourOwnScoreWFlow{}, against \RQfourWFlow{} with the schedule truth. Retrained and scored on their own labels, as a user of one release alone would work, the arms give $W$ of \RQfourOwnRetrWFlow{} at the flow unit and \RQfourOwnRetrWConv{} at the conversation unit (difference \RQfourOwnRetrDiff{}, interval \RQfourOwnRetrDiffLo{} to \RQfourOwnRetrDiffHi{}). H2 still holds, for \RQfourOwnRetrHtwoAbove{} of \RQfourOwnRetrHtwoElig{} eligible classes, but the share direction changes for \RQfourOwnRetrKfive{} of \RQfourClassesOf{} classes. The fall in the exporter share from flow to conversation unit therefore belongs to the setting with one ground truth, and the labels are a second channel through which a release changes scores, which the schedule truth closes by design.

\paragraph{Episodes are not a usable unit on \cicids{}.} An episode groups the attack conversations of one class between one pair of hosts: on \cicids{} within one window of the attack schedule, and on UNSW-NB15, which has no schedule windows, as long as consecutive conversations start at most \setnum{300} seconds after the previous ones end. Episodes were frozen with the instrument and are scored here as a sensitivity unit. On \cicids{} each attack class has at most \RQfourEpisodeMax{} test episodes (\RQfourEpisodeMaxClass{}), and in every seed every attack episode has conversations on both sides of the conversation-grouped split. Claims at the level of an attack campaign cannot be tested on this benchmark, and the conversation split is not independent at that level. On UNSW-NB15 episodes are numerous, but many straddle the split too: \RQfourEpisodeStradDos{} percent of DoS test episodes and \RQfourEpisodeStradExploits{} percent of Exploits test episodes, depending on the seed.

\paragraph{Answer to RQ4.} Conversation-level scoring with conversation-grouped splits does not make model rankings more consistent across exporters. H4 was not supported by the test fixed in advance, with intervals that include zero; with five models the difference on \cicids{} lies below zero, so rankings agree less at the conversation unit, and with three models it reversed when attempted attacks were relabelled benign at test time without retraining. The scoring unit does remove most of the exporter's share of per-class F1 on \cicids{} and raises the share of between-release comparisons that are equivalent within \setnum{0.05}, but it does not make F1 exporter invariant in the sense of Definition~\ref{def:invariance}. What remains on \cicids{} has one identified cause, a short end-of-connection record that the schedule labels as the attack in attack conversations and that makes benign conversations false alarms; a byte-weighted rule or payload-aware labels remove it. On UNSW-NB15 the conversation unit also removes most of the exporter's lead, and the residual that remains in \RQfourLargeUnsw{} classes comes mostly from CIC-UNSW-NB15, through attacks confused with other attack classes, a kind we have not explained.
